\documentclass[10pt,a4paper]{amsart}
\usepackage[margin=19mm]{geometry}
\usepackage{amsmath,amssymb,mathtools}
\usepackage[scr=rsfs,cal=euler]{mathalfa}
\usepackage{xfrac}
\usepackage[unicode,hidelinks,bookmarks=false]{hyperref}
\newcommand{\ie}{i.e.\ }
\newcommand{\cf}{cf.\ }
\newcommand{\resp}{respectively\ }
\newcommand{\Subsection}{\subsection}
\providecommand{\nobreakdash}{\nobreak\hbox{-}}
\setlength{\emergencystretch}{2em}
\multlinegap0pt
\usepackage{bm}
\usepackage{xcolor}
\usepackage{enumerate}
\usepackage{mathtools}
\usepackage{tikz}
\usepackage{amssymb}
\usepackage{algorithm}\usepackage{algpseudocode}
\newtheorem{assumption}{Assumption}
\newtheorem{proposition}{Proposition}
\newcommand{\KL}{D_{\mathrm{KL}}}
\newcommand{\R}{\mathbb{R}}
\newcommand{\map}[3]{#1:#2 \rightarrow #3}
\newcommand{\tr}{\mathrm{tr}}
\title{From Schr\"odinger Bridge to Optimal Transport over Sub-Riemannian Manifolds}
\author{Daniel Owusu Adu, Karthik Elamvazhuthi, Bahman Gharesifard}
\date{Source: arXiv:2605.11429v2 (CC BY 4.0)}
\begin{document}
\maketitle

\noindent\emph{Source and licence.} From Schr\"odinger Bridge to Optimal Transport over Sub-Riemannian Manifolds. Version arXiv:2605.11429v2 (CC BY 4.0), distributed by arXiv under the Creative Commons Attribution 4.0 International licence, http://creativecommons.org/licenses/by/4.0/, as recorded in the arXiv licence metadata for this exact version and shown on the abstract page https://arxiv.org/abs/2605.11429v2. Full licence notice: \texttt{licenses/arxiv-2605.11429-CC-BY-4.0.txt}.\par\smallskip

\noindent\textit{Editorial scope.} Counted unit: the article's proposition attainment (source0.tex lines 617--651). The article's complete original proof of it is delivered with it (source0.tex lines 652--751, one \texttt{proof} environment of 100 lines). No mathematics is written, repaired or re-derived: the statement and the proof are the article's own lines, in the article's own notation, and no retained line is substituted or re-flowed.  Retained uncounted as the source-local context the proof needs: lines 247--252; lines 337--339; lines 351--354; lines 426--428; lines 527--529; lines 558--568.  Nothing was omitted from any retained span, so the package carries no omission record.  No retained line contains a citation, so the wrapper carries no bibliography and no bibliography entry.  No retained line contains a figure environment, an image-inclusion command or drawing code, so no image is delivered and the package declares no asset dependency.  Editorial formatting only, in addition: an AMS \texttt{amsart} preamble that restates the article's own declarations and macros used by the retained lines, namely the environments \texttt{assumption}, \texttt{proposition} on separate counters, together with the standard packages those lines need.  The retained environments print in this document as Assumption 1, Proposition 1 in order of appearance, whereas the article prints the same items with the numbering of its own full document.  The complete native source of the article as distributed in the arXiv e-print for this exact version is bundled unchanged as \texttt{sources/200356/source0.tex}.
\begin{assumption}\label{ass:Controllability condition}
The vector fields ${g_1, \dots, g_m}$ satisfy the bracket-generating condition. Specifically, we assume that the Lie algebra generated by $\{g_1, \dots, g_m\}$ spans the tangent space $T_y\mathbb{R}^d$ at every point $y \in \mathbb{R}^d$:
\[
\operatorname{Lie}(g_1, \dots, g_m)(y) = \mathbb{R}^d \quad \text{ for all $y \in \mathbb{R}^d$.}
\]
\end{assumption}

\begin{equation}\label{kl_control_cost}
\epsilon\KL(\mathrm{P}^u \| \mathrm{R}^{\epsilon,\mu_0}) = \mathbb{E}_{\mathrm{P}^u} \left[ \frac12 \int_0^{t_f} \|u(t,  X_t)\|^2 dt \right],
\end{equation}

\begin{equation}\label{eq:Ito-generator}
\mathcal L_\epsilon 
= b_\epsilon\cdot\nabla_x +\frac{\epsilon}{2}\mathrm{tr}\!\big(G\nabla_x^2\big)
\end{equation}

\begin{equation}\label{eq:Sch-multplication_system}
\phi_0(x)(\mathscr P_{0,t_f} \phi_f)(x)=\rho_0(x)\qquad\text{and}\qquad (\mathscr P_{0,t_f}^* \phi_0)(y)\phi_f(y)=\rho_f(y),
\end{equation}

\begin{multline}\label{eq: finite_energy_bound}
\min_{\mathrm{P}^u \in \mathcal C_{\mathrm{SB}}} \KL(\mathrm{P}^u \| \mathrm{R}^{\epsilon,\mu_0})\leq \KL(\mu_0\otimes\mu_f \| r)\\=\int_{\R^d}\rho_f(y)\log(\rho_f(y))dy-\int_{\R^d\times\R^d}\log(p_{t_f,\epsilon}(x,y))\rho_0(x)\rho_f(y)dxdy
\end{multline}

\begin{align}
   \partial_t \rho_{\epsilon}
+ \nabla_x \cdot \left(b_{\epsilon}\rho_{\epsilon}+gm_{\epsilon}\right)
- \frac{\epsilon}{2} \sum_{i,j=1}^d 
\partial_{x_i x_j}^2 \Big(G_{ij} \rho_{\epsilon} \Big)
&= 0\quad\text{in }\quad \bigl(C_c^\infty(Q)\bigr)^{*}\label{eq: fb_constraint}\\
    \rho_{\epsilon}(0, \cdot) = \rho_0, 
    \qquad 
    \rho_{\epsilon}(t_f, \cdot) &= \rho_{t_f},\quad\quad\text{in\quad $\mathcal P(\R^d)$.} 
    \label{eq:endpoint_constraints}
\end{align}

\begin{proposition}\label{prop: Strong duality and attainment}
Let 
\begin{equation}\label{eq: dual_set_thm}
 \Lambda := \Bigg\{ \lambda \in C^{1,2}([0,t_f]\times\mathbb{R}^d) \;\Bigg|\; \partial_t\lambda + b_\epsilon\cdot\nabla\lambda + \frac12\langle \nabla\lambda, G\nabla\lambda\rangle + \frac{\epsilon}{2}\operatorname{tr}(G\nabla^2\lambda) \le 0 \quad \text{on } (0,t_f)\times\mathbb{R}^d \Bigg\}   
\end{equation}
and define the dual functional $\map{\mathcal D}{\Lambda}{\{\R\cup-\infty\}}$ as
\begin{equation}\label{eq:dual-func-thm}
\mathcal D(\lambda):=\int_{\R^d}\lambda(t_f,x)\rho_f(x) dx
-\int_{\R^d}\lambda(0,x)\rho_0(x) dx.
\end{equation}

Assume that there exist strictly positive functions $(\varphi_\epsilon,\hat\varphi_\epsilon)$ that solve: %the Schr\"odinger system
\begin{subequations}\label{eq:Schr-system-thm}
\begin{align}
\partial_t\varphi_\epsilon + \mathcal L_\epsilon \varphi_\epsilon &=0,\label{eq:Schr-fwd}\\
\partial_t\hat\varphi_\epsilon - \mathcal L_\epsilon^* \hat\varphi_\epsilon &=0,\label{eq:Schr-bwd}\\
\varphi_\epsilon(0,\cdot)\hat\varphi_\epsilon(0,\cdot)=\rho_0,\qquad
\varphi_\epsilon(t_f,\cdot)\hat\varphi_\epsilon(t_f,\cdot)&=\rho_f,\label{eq:Schr-bc}
\end{align}
\end{subequations}
where $\mathcal L_\epsilon$ is in~\eqref{eq:Ito-generator}. Then strong duality holds: 
\[ 
\inf_{\substack{(\rho_{\epsilon}, m_{\epsilon}) \\ \text{\eqref{eq: fb_constraint}-\eqref{eq:endpoint_constraints}}}} 
\mathcal{J}(\rho_{\epsilon}, m_{\epsilon}) = \sup_{\lambda_{\epsilon}\in\Lambda}\mathcal D(\lambda_{\epsilon})
\]
and the supremum is attained at
\begin{equation}\label{eq:optimal-dual}
\lambda_\epsilon^\star:=\epsilon\log\varphi_\epsilon.
\end{equation}
Moreover, the primal infimum is uniquely attained at
\begin{align}\label{eq:optimal-primal}
&\rho_\epsilon^\star:=\varphi_\epsilon\hat\varphi_\epsilon\in C^\infty((0,t_f)\times\R^d),\nonumber\\\
&m_\epsilon^\star:=\rho_\epsilon^\star g^\top\nabla\lambda_\epsilon^\star = \epsilon \hat\varphi_\epsilon g^\top\nabla\varphi_\epsilon\in C^\infty((0,t_f)\times\R^d).
\end{align}    
\end{proposition}

\begin{proof}
We establish the result in three steps: weak duality, feasibility of the constructed pair, and equality verifying attainment.

For weak duality, for any primal feasible $(\rho,m)$ and dual feasible $\lambda\in\Lambda$, we test the Fokker-Planck constraint in~\eqref{eq: fb_constraint}-\eqref{eq:endpoint_constraints} against $\lambda$ (approximating by compactly supported test functions if necessary). Hence integration by parts yields
\begin{align*}
\mathcal D(\lambda_{\epsilon})
&=\int_0^{t_f}\!\!\int_{\R^d}
\Bigl[\rho_{\epsilon}\bigl(\partial_t\lambda_{\epsilon}+b_\epsilon\cdot\nabla\lambda_{\epsilon}
+\tfrac{\epsilon}{2}\mathrm{tr}(G\nabla^2\lambda_{\epsilon})\bigr)
+ m\cdot(g^\top\nabla\lambda_{\epsilon})\Bigr] dx dt.
\end{align*}
Since $\lambda_{\epsilon}\in\Lambda$ in~\eqref{eq: dual_set_thm}, we get that
\[
\partial_t\lambda_{\epsilon}+b_\epsilon\cdot\nabla\lambda_{\epsilon}
+\tfrac{\epsilon}{2}\mathrm{tr}(G\nabla^2\lambda_{\epsilon})
\le -\tfrac12\langle\nabla\lambda_{\epsilon},G\nabla\lambda_{\epsilon}\rangle
=-\tfrac12\|g^\top\nabla\lambda_{\epsilon}\|^2,
\]
and hence
\[
\mathcal D(\lambda_{\epsilon})\le \int_0^{t_f}\!\!\int_{\R^d}
\Bigl(m_{\epsilon}\cdot(g^\top\nabla\lambda_{\epsilon})-\tfrac{\rho_{\epsilon}}{2}\|g^\top\nabla\lambda_{\epsilon}\|^2\Bigr) dx dt.
\]
Applying Young's inequality pointwise gives
\begin{equation}\label{eq: Young's inequality}
m_{\epsilon}\cdot(g^\top\nabla\lambda_{\epsilon})=\left(\frac{m_{\epsilon}}{\sqrt{\rho_{\epsilon}}}\right)\left(\sqrt{\rho_{\epsilon}}(g^\top\nabla\lambda_{\epsilon}\right)
\le \frac{\|m_{\epsilon}\|^2}{2\rho_{\epsilon}}+\frac{\rho_{\epsilon}}{2}\|g^\top\nabla\lambda_{\epsilon}\|^2,
\end{equation}
and hence we obtain 
\[
\mathcal D(\lambda_{\epsilon})\le \mathcal J(\rho_{\epsilon},m_{\epsilon})
\]
for all $\lambda_{\epsilon}\in\Lambda$ in~\eqref{eq: dual_set_thm} and the pairs $(\rho_{\epsilon},m_{\epsilon})$ in~\eqref{eq: fb_constraint}-\eqref{eq:endpoint_constraints}. Thus 
\begin{equation}\label{eq: weak duality} 
\inf_{\substack{(\boldsymbol\mu_{\epsilon}, \mathbf m_{\epsilon})\in \mathcal C \\ \text{\eqref{eq: fb_constraint}-\eqref{eq:endpoint_constraints}}}} 
\mathcal{J}(\rho_{\epsilon},m_{\epsilon}) \geq \sup_{\lambda_{\epsilon}\in\Lambda}\mathcal D(\lambda_{\epsilon}).
\end{equation}
%\margin{should $ \rho_{\epsilon}^\star$ be $\rho_\epsilon^\star$}
For feasibility and equality in the HJB, we consider the function tuple $(\lambda_{\epsilon}^\star,\rho_{\epsilon}^\star,m_{\epsilon}^\star)$ in \eqref{eq:optimal-dual}--\eqref{eq:optimal-primal}. From \eqref{eq:Schr-bc}, $\rho_{\epsilon}^\star$ satisfies the endpoint conditions~\eqref{eq:endpoint_constraints}. A direct computation using $\lambda_{\epsilon}^\star=\epsilon\log\varphi_{\epsilon}$ and \eqref{eq:Schr-fwd} gives
\begin{multline*}
\partial_t\lambda^\star+b_\epsilon\cdot\nabla\lambda^\star
+\tfrac12\langle\nabla\lambda^\star,G\nabla\lambda^\star\rangle
+\tfrac{\epsilon}{2}\tr(G\nabla^2\lambda^\star) \\
= \epsilon\frac{\partial_t\varphi_{\epsilon}+b_\epsilon\cdot\nabla\varphi_{\epsilon}+\tfrac{\epsilon}{2}\tr(G\nabla^2\varphi_{\epsilon})}{\varphi_{\epsilon}}
-\tfrac{\epsilon^2}{2}\frac{\langle\nabla\varphi_{\epsilon},G\nabla\varphi_{\epsilon}\rangle}{\varphi_{\epsilon}^2}
+\tfrac{\epsilon^2}{2}\frac{\langle\nabla\varphi_{\epsilon},G\nabla\varphi_{\epsilon}\rangle}{\varphi_{\epsilon}^2}
\\= \epsilon\frac{\partial_t\varphi_{\epsilon}+\mathcal L_\epsilon\varphi_{\epsilon}}{\varphi_{\epsilon}}=0.
\end{multline*}
Thus $\lambda^\star\in\Lambda$, achieving the equality in \eqref{eq: dual_set_thm}. To verify primal feasibility, note that $m_{\epsilon}^\star=\rho_{\epsilon}^\star g^\top\nabla\lambda^\star$ by construction. Using \eqref{eq:Schr-system-thm}, the pair $(\rho_{\epsilon}^\star,m_{\epsilon}^\star)$ satisfies
\begin{align*}
\partial_t\rho_{\epsilon}^\star
&= \hat\varphi_{\epsilon}\,\partial_t\varphi_{\epsilon} + \varphi_{\epsilon}\,\partial_t\hat\varphi_{\epsilon}
= -\hat\varphi_{\epsilon}\,\mathcal L_\epsilon\varphi_{\epsilon} + \varphi_{\epsilon}\,\mathcal L_\epsilon^\ast\hat\varphi_{\epsilon} \\
&= -\hat\varphi_{\epsilon}\Big(b_\epsilon\cdot\nabla\varphi_{\epsilon} + \frac{\epsilon}{2}\tr(G\nabla^2\varphi_{\epsilon})\Big)
\;+\;\varphi_{\epsilon}\Big(-\nabla\!\cdot(b_\epsilon\hat\varphi_{\epsilon})+\frac{\epsilon}{2}\sum_{i,j=1}^d\partial_{ij}(G_{ij}\hat\varphi_{\epsilon})\Big) \\
&= -\nabla\!\cdot(b_\epsilon\,\rho_{\epsilon}^\star)
\;+\;\frac{\epsilon}{2}\sum_{i,j=1}^d \partial_{ij}\!\big(G_{ij}\rho_{\epsilon}^\star\big)
\;-\;\epsilon\,\nabla\!\cdot\!\big(\hat\varphi_{\epsilon}\,G\nabla\varphi_{\epsilon}\big).
\end{align*}
where the cross terms combine to yield precisely $\nabla\cdot(gm_{\epsilon}^\star) = \nabla\cdot(\epsilon\hat\varphi_{\epsilon} G\nabla\varphi_{\epsilon})$. Hence the pair $(\rho_{\epsilon}^\star,m_{\epsilon}^\star)$ solves~\eqref{eq: fb_constraint}. Furthermore, since
%\[
%\frac{\|m^*\|^2}{\rho_{\epsilon}^\star} = \frac{\epsilon^2 \hat{\varphi}^2 \|g^\top \nabla\varphi\|^2}{\varphi\hat{\varphi}} = \epsilon^2 \hat{\varphi} \frac{\|g^\top \nabla\varphi\|^2}{\varphi}
%\]
%and hence
%Using the identity $\partial_t\varphi + \mathcal{L}_\epsilon\varphi = 0$, one shows that
\[
\int_Q \frac{\|m_{\epsilon}^\star\|^2}{2\rho_{\epsilon}^\star} dxdt = \epsilon^2 \int_Q \hat{\varphi_{\epsilon}} \frac{\|g^\top \nabla\varphi_{\epsilon}\|^2}{2\varphi_{\epsilon}} dxdt= \epsilon^2\int_Q \hat{\varphi_{\epsilon}}  \varphi_{\epsilon} \frac{\|g^\top \nabla(\log\varphi_{\epsilon})\|^2}{2}   dxdt,
\]
%\margin{Where is the $1/2 $ coming from? I think there is a minor error here}
and
\[
\epsilon^2\int_Q \hat{\varphi_{\epsilon}}   \varphi_{\epsilon}   \frac{\|g^\top \nabla(\log\varphi_{\epsilon})\|^2}{2}   dxdt=\int_Q  \frac12\|u^*\|^2\rho_{\epsilon}^\star   dxdt
\]
from~\eqref{kl_control_cost} and the energy estimates in~\eqref{eq: finite_energy_bound} we conclude that
\[
\int_Q \frac{\|m_{\epsilon}^\star\|^2}{\rho_{\epsilon}^\star} dxdt<\infty.
\]

For attainment and strong duality, note that if $m_{\epsilon}^\star=\rho_{\epsilon}^\star g^\top\nabla\lambda_{\epsilon}^\star$, then the equality holds in Young's inequality in~\eqref{eq: Young's inequality} and equality holds in~\eqref{eq: weak duality}. Therefore
\[
\mathcal D(\lambda_{\epsilon}^\star) = \mathcal J(\rho_{\epsilon}^\star,m_{\epsilon}^\star).
\]
Combined with weak duality, this yields
\[
\mathcal D(\lambda^\star) = \mathcal J(\rho_{\epsilon}^\star,m_{\epsilon}^\star) = \inf\mathcal J = \sup_{\Lambda}\mathcal D,
\]
proving that $\lambda_{\epsilon}^\star$ attains the dual supremum, $(\rho_{\epsilon}^\star,m_{\epsilon}^\star)$ attains the primal infimum, and there is no duality gap.

For regularity, note that the solution for~\eqref{eq:Schr-fwd}-\eqref{eq:Schr-bc} is characterized as
\[
\varphi_\epsilon(t,\cdot) = e^{(t_f-t)\mathcal{L}_\epsilon}\varphi_{t_f}, \qquad 
\hat{\varphi}_\epsilon(t,\cdot) = e^{t\mathcal{L}_\epsilon^*}\hat{\varphi}_0,
\]
with $\hat{\varphi}_0 = \rho_0 / \varphi(0,\cdot)$ and $\varphi_{t_f}=\phi_f$ in~\eqref{eq:Sch-multplication_system},  where the forward and backward semigroups $P_t = e^{t\mathcal{L}_\epsilon^*}$ and $
Q_t = e^{t\mathcal{L}_\epsilon}$ are defined through the actions
\begin{equation}\label{eq: action_semigroup}
(P_t \varphi_0)(y) = \int_{\R^d} p_{t,\epsilon}(x,y) \varphi_0(x)   dx\quad\text{and}\quad (Q_t \varphi_f)(x) = \int_{\R^d} p_{t,\epsilon}(x,z) \varphi_f(z)   dz    
\end{equation}
respectively and $p_{t,\epsilon}(x,y)$ is the transition density function. Under Assumption~\ref{ass:Controllability condition}, since the kernels $p_{t,\epsilon}(x,y)$ are smooth and strictly positive for all $t>0$ we have that the functions $\varphi$ and $\hat{\varphi}$  inherit this smoothness and strict positivity. Hence $\rho_\epsilon^* = \varphi_\epsilon\hat{\varphi}_\epsilon \in C^\infty((0,t_f)\times\R^d)$ and $\rho_{\epsilon}^\star > 0$ everywhere. Moreover, $m_{\epsilon}^\star$ is smooth by the smoothness of $\varphi_\epsilon$, $\hat{\varphi}_\epsilon$, and $g$. The uniqueness follows from the strict convexity of $\mathcal{J}$ and the linearity of the constraint set.
\end{proof}

\end{document}
